\begin{statementbox}
	\textbf{\textcolor{CStatement}{条件}}
	\begin{description}[style=nextline, leftmargin=2.8em, labelsep=0.5em, itemsep=0.35em]
		\item[\textbf{\textcolor{CStatement}{条件1（对数定义域）：}}] $\ln x \le x-1$ 要求 $x>0$。
		\item[\textbf{\textcolor{CStatement}{条件2（指数定义域）：}}] $e^x \ge x+1$ 和 $e^x \ge ex$ 对任意实数 $x$ 成立。
	\end{description}

	\textbf{\textcolor{CStatement}{结论}}
	\begin{description}[style=nextline, leftmargin=2.8em, labelsep=0.5em, itemsep=0.45em]
		\item[\textbf{\textcolor{CStatement}{结论一（对数放缩）：}}]
			\textbf{\textcolor{CStatement}{直观描述：}}
			对数函数 $\ln x$ 在 $x>0$ 时被一次函数 $x-1$ 从上方控制。
			\[
				\ln x \le x-1
			\]
			等号当且仅当 $x=1$。

		\item[\textbf{\textcolor{CStatement}{结论二（指数放缩一）：}}]
			\textbf{\textcolor{CStatement}{直观描述：}}
			指数函数 $e^x$ 在全体实数上被一次函数 $x+1$ 从下方控制。
			\[
				e^x \ge x+1
			\]
			等号当且仅当 $x=0$。

		\item[\textbf{\textcolor{CStatement}{结论三（指数放缩二）：}}]
			\textbf{\textcolor{CStatement}{直观描述：}}
			指数函数 $e^x$ 被过原点的切线 $ex$ 从下方控制。
			\[
				e^x \ge ex
			\]
			等号当且仅当 $x=1$。

		\item[\textbf{\textcolor{CStatement}{推广结论（组合放缩）：}}]
			\textbf{\textcolor{CStatement}{直观描述：}}
			将对数放缩与指数放缩一相加，消去一次项得到常数下界。
			\[
				e^x - \ln x \ge 2 \quad (x>0)
			\]
			注意此处等号无法同时取到，实际恒大于 $2$。
	\end{description}
\end{statementbox}
